
\subsection{Core Principles}\label{sec:core_principles}
Building on the two-stage framework, we conclude four core principles that define a well-designed counterfactual query system.

\paragraph{\textbf{P1 (Traceability).}} Counterfactual responses must remain interpretable in the original domain. For any $x' \in S^\star$, the system outputs an explicit \emph{domain-level edit description} $\Delta(x_q,x')$, such as changed attributes with before--after values in tabular data or a valid sequence of edit operations in structured modalities. The representation and indexing mechanism must preserve sufficient information to reconstruct $\Delta(x_q,x')$ without ambiguity. The transformation cost induced by $\omega$, i.e., $\omega(x_q,x')$, must be grounded in this edit and admit a mapping back to $\Delta(x_q,x')$, ensuring that optimization remains tied to domain-level transformations rather than abstract similarity.
\paragraph{\textbf{P2 (Cost faithfulness).}} The transformation cost induced by $\omega$, i.e., $\omega(x_q,x)$, encodes the semantic notion of transformation effort and must align with domain actionability: if $\omega(x_q,x_a) < \omega(x_q,x_b)$, then reaching $x_a$ from $x_q$ should not require more effort than reaching $x_b$ under the intended intervention model. Infeasible transformations are assigned infinite cost, ensuring that admissibility and effort are jointly captured within $\omega$.
\paragraph{\textbf{P3 (Efficiency).}} The system must support scalable and predictable access to admissible candidates $\{x \in \mathcal D : \omega(x_q,x) < \infty\}$, including bounded resource usage, stable performance under typical workloads, and controllable approximation behavior when trading accuracy for latency.
\paragraph{\textbf{P4 (Collective Validity).}} Beyond individual admissibility, the returned set must optimize the user-specified objective $\Omega$. For a query $(q,\omega,\Omega,k)$, the response set $S^\star$ satisfies $|S^\star|=k$ and maximizes $\Omega(S;q)$ over admissible candidates. Because $\Omega$ captures collective properties such as diversity, coverage, or fairness, optimizing $\Omega$ may trade off against strict per-instance minimality under $\omega(x_q,x)$, and this trade-off should be explicit and controllable.

\smallskip\noindent To this end, these four principles jointly characterize a counterfactual query system at the representation, semantic, retrieval, and set-selection levels, distinguishing structural correctness from set-level utility. Sections~\ref{sec:instance} and~\ref{sec:set} describe how $\omega$ and $\Omega$ are instantiated in practice and how they shape system behavior across diverse data modalities and applications.